\section{Introduction}

Reinforcement learning (RL) depends critically on reward functions that express task intent while providing informative feedback for policy optimization~\cite{sutton2018rl}. Reward shaping can improve credit assignment, but only restricted transformations are guaranteed to preserve the optimal policy~\cite{ng1999policy}. Inverse reward design further shows that a written reward should be treated as imperfect evidence of designer intent rather than as the intent itself~\cite{hadfield2017inverse}. Reward misspecification can therefore produce unsafe or unintended behavior even when the specified signal appears semantically reasonable~\cite{amodei2016concrete}.

Reward engineering is consequently an iterative debugging process rather than a one-shot coding task. Each completed training run produces more than a scalar return: episode length, termination patterns, reward-component activation, magnitude balance, and temporal trends can reveal how the reward program shaped the learned policy. Converting these heterogeneous records into a precise reward modification, however, still requires substantial human analysis.

Code-capable large language models (LLMs) can generate reward programs from task descriptions~\cite{language2rewards,text2reward} and iteratively refine them through training feedback and evolutionary search~\cite{eureka,card,revolve,cotreward}. Diagnostic-driven refinement further treats reward failure as a debugging problem~\cite{wang2026diagnostic}. In parallel, LLMs are increasingly organized as agents that observe, reason, and act across trials, using stored feedback to improve later decisions~\cite{react,reflexion,selfrefine}. These developments motivate an agent-oriented view of reward engineering in which policy training supplies observations and reward-code revision constitutes an external action.

This paper explores an agent-oriented view of reward engineering through CREATE, a \emph{Closed-Loop Reward Engineering Agent with Training Evidence}. CREATE places a reward-engineering agent above an inner RL policy learner. After each policy-training run, the reward-engineering agent uses explicit tool calls to inspect training feedback, reward and lineage history, and component-level trends while it reasons about the next repair. It then diagnoses one principal semantic failure and applies a localized edit: L1 tunes parameters, L2 refactors one semantic, and L3 redesigns the reward structure after repeated local failure. Intervention memory retains this evidence chain across rounds, while a separate best archive protects the strongest solution from later exploratory regressions.

Experiments on LunarLander-v3 compare CREATE against budget-matched independent generation and baseline, with two mechanism ablations isolating evidence enrichment and hierarchical semantic editing. BipedalWalker-v3 provides a secondary cross-task validation. Under a ten-evaluation budget, the experiments demonstrate that structured observation, persistent state, and bounded reward actions convert failed policy training into reliable reward repair.
